
\chapter*{ABSTRACT}

\footnotesize{
\begin{description}
  \item[TITLE:] \MakeUppercase{Detection of Diabetic Retinopathy in Fundus Images Using Deep Learning Algorithms.}
  \astfootnote{Bachelor Thesis}
  \item[AUTHOR:] {\autora}, {\autorb}
  \asttfootnote{\facultad. \escuela. Advisor: \director. Co-advisor: \codirectora}
  \item[KEYWORDS:] DIABETIC RETINOPATHY, CONVOLUTIONAL NEURONAL NETWORKS, DEEP LEARNING, FUNDUS IMAGES, DIGITAL IMAGE PROCESSING, TRANSFER LEARNING.
  \item[DESCRIPTION:]\hfill \\  Diabetic retinopathy (DR) is an ocular complication caused by damage to the blood vessels in the retina due to elevated blood sugar levels, commonly found in people with diabetes. It is considered one of the leading causes of blindness worldwide, highlighting the importance of early diagnosis to prevent or minimize visual impairment. This disease is typically diagnosed through eye exams that assess the retina. However, manual diagnosis is challenging because it is costly and time-consuming, limiting access in areas with scarce medical resources. This underscores the need for automated methods that enable faster and more accurate disease detection. This study proposes two models based on convolutional neural networks (CNNs) that not only detect the disease but also classify its different severity levels. The model is trained using a public dataset and later tested with a private dataset provided by the Santander Ophthalmological Foundation. This study addresses the issue of data imbalance through preprocessing and image augmentation techniques. Additionally, various architectures are built using pretrained networks to extract features from the images. The models performance is evaluated and compared using different metrics, where the best binary classification model achieves a sensitivity of 96.5\%, and the best multiclass classification model reaches 76.8\%

\end{description}}\normalsize
% ------------------------------------------------------------------------ 
